
% ────────────────────────────────────────────────────────
\section{O Integral Definido}

% ────────────────────────────────────────────────────────
\subsection{Enquadramento Teórico}

\begin{defbox}[title={Integral de Riemann}]
  Seja $f:[a,b]\to\mathbb{R}$ limitada. Para uma partição
  $P=\{a=x_0<x_1<\cdots<x_n=b\}$, a \emph{soma de Riemann} é
  \[
    S(f,P) = \sum_{i=1}^n f(c_i)(x_i - x_{i-1}),\quad c_i\in[x_{i-1},x_i].
  \]
  Se $S(f,P)$ convergir para um limite $L$ quando a norma da partição
  $\|P\|\to 0$ (independentemente da escolha de $c_i$), então $f$ é
  \emph{integrável no sentido de Riemann} e $\displaystyle\int_a^b f(x)\,dx = L$.
\end{defbox}

\begin{thmbox}[title={Propriedades Principais do Integral Definido}]
\begin{align*}
  &\int_a^b [\alpha f+\beta g]\,dx = \alpha\int_a^b f\,dx + \beta\int_a^b g\,dx
    \quad (\text{linearidade})\\[4pt]
  &\int_a^b f\,dx = \int_a^c f\,dx + \int_c^b f\,dx
    \quad (\text{aditividade em } [a,c]\cup[c,b])\\[4pt]
  &f(x)\geq 0 \implies \int_a^b f\,dx \geq 0
    \quad (\text{positividade})
\end{align*}
\end{thmbox}
